\clearpage\thispagestyle{empty}
\pdfbookmark[0]{Quellen zu Bedeckung von 42 Omega Strich Geminorum (1841)}{occultation1841-sources}
\begin{center}\textbf{Quellen zu Bedeckung von 42 Omega Strich Geminorum (1841)}\end{center}
\noindent Diese KI-Übertragung ins heutige Deutsch folgt dem Text auf S.~464 der \href{https://gdz.sub.uni-goettingen.de/id/PPN236007467}{\textit{Werke}, Band~VI} (Göttingen, 1874). Die Mitteilung von Heinrich Christian Schumacher gibt Beobachtungen von Carl Friedrich Gauß und Dr. Goldschmidt vom 23. Mai 1841 wieder. Das Heft des Erstdrucks ist auf den 5. August 1841 datiert. \textit{Astronomische Nachrichten}, Band~XVIII, Nr.~430, \href{https://api.digitale-sammlungen.de/iiif/image/v2/bsb10049443_00194/full/1800,/0/default.jpg}{Sp.~367--368}.
\par\smallskip\noindent Die Werkausgabe nennt das Jahr und gibt Schumachers Unterschrift wieder. Der Erstdruck hat eine eigene Überschrift, \textit{Sternbedeckung}, und verwendet andere Zeichen für die Zeiteinheiten. Die Übertragung folgt der Werkausgabe. KI-Übertragung und KI-Satz.
\clearpage
